% Part:sets-functions-relations
% Chapter: size-of-sets
% Section: equinumerous-sets

\documentclass[../../../include/open-logic-section]{subfiles}

\begin{document}

\olfileid[fa-IR]{sfr}{siz}{equ}

\olsection{هم‌توانی}

تصوری شهودی از «اندازه» مجموعه‌ها داریم که برای مجموعه‌های متناهی به‌خوبی
کار می‌کند. اما مجموعه‌های نامتناهی چه؟ اگر بخواهیم روشی صوری برای
مقایسهٔ اندازهٔ دو مجموعه، با \emph{هر} اندازه‌ای که داشته باشند، به
دست دهیم، بهتر
است کار را با تعریف اینکه چه هنگام دو مجموعه هم‌اندازه‌اند آغاز کنیم.
فرگه چنین می‌گوید:
\begin{quote}
  اگر پیشخدمتی بخواهد مطمئن شود که روی میز دقیقاً به تعداد بشقاب‌ها
  چاقو چیده است، لازم نیست هیچ‌یک را بشمارد؛ کافی است در سمت راست هر
  بشقاب یک چاقو بگذارد، به‌گونه‌ای که هر چاقوی روی میز در سمت راست
  یکی از بشقاب‌ها قرار گیرد. به این ترتیب، بشقاب‌ها و چاقوها به‌طور
  یکتا با یکدیگر متناظر می‌شوند، آن هم به‌واسطهٔ همان رابطهٔ مکانی.
  \citep[\S70]{Frege1884}
\end{quote}
بینش این قطعه را می‌توان با تعریفی صوری آشکار کرد:

\begin{defn}\ollabel{comparisondef}
  مجموعهٔ $A$ با $B$ \emph{هم‌توان} است، که آن را با $\cardeq{A}{B}$
  می‌نویسیم، اگر و تنها اگر !!a{bijection} به‌صورت
  $f \colon A \to B$ وجود داشته باشد.
\end{defn}

\begin{prop}\ollabel{equinumerosityisequi}
هم‌توانی یک رابطهٔ هم‌ارزی است.
\end{prop}

\begin{proof}
باید نشان دهیم که هم‌توانی بازتابی، متقارن و تراگذری است. فرض کنید
$A, B$ و $C$ مجموعه باشند.

\emph{بازتابی‌بودن.} نگاشت همانیِ $\Id{A} \colon A \to A$، که در آن
$\Id{A} (x) = x$ برای هر $x \in A$، !!a{bijection} است. پس
$\cardeq{A}{A}$.

\emph{تقارن.} فرض کنید $\cardeq{A}{B}$؛ یعنی !!a{bijection} به‌صورت
$f\colon A \to B$ وجود دارد. چون $f$ !!{bijective} است، وارون آن
$f^{-1}$ وجود دارد و آن نیز !!{bijective} است. بنابراین
$f^{-1}\colon B \to A$ !!a{bijection} است؛ پس $\cardeq{B}{A}$.

\emph{تراگذری.} فرض کنید $\cardeq{A}{B}$ و $\cardeq{B}{C}$؛ یعنی
!!{bijection}s زیر وجود دارند: $f\colon A \to B$ و $g\colon B \to
C$. آنگاه ترکیب $\comp{f}{g}\colon A \to C$ !!{bijective}
است؛ پس $\cardeq{A}{C}$.
\end{proof}

\begin{prop}
اگر $\cardeq{A}{B}$، آنگاه $A$ !!{enumerable} است اگر
و تنها اگر $B$ چنین باشد.
\end{prop}

\begin{editorial}
برهان زیر از \olref[enm]{defn:enumerable} استفاده می‌کند، اگر
\olref[enm]{sec} گنجانده شده باشد، و در غیر این صورت از
\olref[enm-alt]{defn:enumerable}.
\end{editorial}

\begin{proof}
فرض کنید $\cardeq{A}{B}$؛ پس یک !!{bijection} به‌صورت $f \colon A
\to B$ وجود دارد، و فرض کنید $A$ !!{enumerable} است.
\oliflabeldef{sfr:siz:enm:defn:enumerable}{
  آنگاه یا $A = \emptyset$، یا تابعی !!a{surjective} به‌صورت
  $g\colon \PosInt \to A$ وجود دارد. اگر $A = \emptyset$، آنگاه $B = \emptyset$
  نیز (زیرا در غیر این صورت !!a{element}~$y \in B$ وجود می‌داشت،
  اما هیچ عضوی مانند $x \in
  A$ با $g(x) = y$ وجود نداشت). از سوی دیگر، اگر $g\colon \PosInt \to A$
  !!{surjective} باشد، آنگاه $\comp{g}{f} \colon \PosInt \to B$
  !!{surjective} است. برای دیدن این مطلب، $y \in B$ را در نظر بگیرید.
  چون $f$ !!{surjective} است، عضوی مانند $x \in A$ وجود دارد که $f(x) = y$.
  چون $g$ !!{surjective} است، عددی مانند $n \in \PosInt$ وجود دارد که $g(n) =
  x$. بنابراین
\[
(\comp{g}{f})(n) = f(g(n)) = f(x) = y
\]
و در نتیجه $\comp{g}{f}$ !!{surjective} است. پس $\comp{g}{f}$
برشماری‌ای از~$B$ است و بنابراین $B$~!!{enumerable} است.}
{
  آنگاه یا $A  = \emptyset$، یا یک !!a{bijection}~$g$ وجود دارد که
  بردش $A$ و دامنه‌اش یا $\Nat$ یا قطعه‌ای آغازین از اعداد طبیعی
  است. اگر $A = \emptyset$، آنگاه $B =
  \emptyset$ نیز (زیرا در غیر این صورت عضوی مانند~$y \in B$ وجود می‌داشت، اما
  هیچ عضوی مانند $x
  \in A$ نبود که $g(x) = y$). پس فرض کنید !!{bijection}~$g$ را
  داریم. آنگاه $\comp{g}{f}$ !!a{bijection} است که بردش~$B$ و دامنه‌اش
  همان دامنهٔ~$g$ است (یعنی یا $\Nat$ یا قطعه‌ای آغازین از آن)؛ پس
  $B$ !!{enumerable} است.}

اگر $B$ !!{enumerable} باشد، نتیجه می‌گیریم که $A$ !!{enumerable} است؛
کافی است استدلال را با !!{bijection} $f^{-1}\colon B \to A$ به‌جای~$f$
تکرار کنیم.
\end{proof}

\begin{prob}
نشان دهید اگر $\cardeq{A}{C}$ و $\cardeq{B}{D}$ و $A \cap B =
C \cap D = \emptyset$، آنگاه $\cardeq{A \cup B}{C \cup D}$.
\end{prob}

\begin{prob}
نشان دهید اگر $A$ نامتناهی و !!{enumerable} باشد، آنگاه
$\cardeq{A}{\Nat}$.
\end{prob}

\end{document}
